\section{Conditional Reflection}

Looking at \eqref{eq:analytical_reflection} one can see that the behavior of the reflection amplitude at zero detuning ($\Delta_a=\Delta_c=0$) is:

\begin{equation}
    r(\Delta = 0) = \mu_{fr} - \mu_{fc}^2\frac{2\kappa_{oc}\gamma}{\kappa\gamma + g^2}
    \label{eq:analytical_reflection_zero_detuning}
\end{equation}

In the event of an uncoupled atom-cavity system ($g=0$), which would be achieved by preparing the atom in the $\ket{u}$ state, the reflection coefficient $r$ is going to be smaller than 0, which would imply a relative phase shift of $\pi$ between the reflected and incoming light field. Placing the atom in the coupled state $\ket{c}$ however, the reflection coefficient $r$ is going to be bigger than 0 (at sufficiently high coopoerativity). Equipped with this knowledge, that there is a relative phase shift of the output field depending on in which state the atom is, allows the realisation of quantum atom-photon gate with the atom being the controlling qubit and the photon being the target qubit.\cite{DuanKimble2004}

\subsection{Reflection in CPF basis}

Changing the nomenclature to match that regularly used in quantum computing, the uncoupled state of the atom is going to be the $\ket{0}$ state, and the coupled one the $\ket{1}$ state. Our photon which as stated before, is going to be either $\pi\text{-polarized}$ or V-polarized, will thus be represented as the input states $\ket{\pi}$ and $\ket{V}$. For the case of our atom being in the $\ket{0}$ state and us sending in a photon which is resonant to the $\pi\text{-mode}$ of the cavity (here: $\ket{\pi}$), the output-field will receive a phase flip, an none for all the other cases either due to being far detuned from the cavity resonance (case of sending $\ket{V}$) and/or atom being in the coupled case $\ket{1}$ and the photon thus getting reflected off of the first cavity mirror:

\begin{align*}
    \ket{0,\pi}&\rightarrow-\ket{0,\pi}\\
    \ket{0,V}&\rightarrow\ket{0,V}\\
    \ket{1,\pi}&\rightarrow\ket{1,\pi}\\
    \ket{1,V}&\rightarrow\ket{1,V}\\
\end{align*}

From this we deduce that output from the atom-photon reflection gives us a CPF gate. This conditional phase shift allows the construction of a universal quantum gate that can be transformed into any two-qubit gate using rotations of the individual qubits, which in the case of the photon can be done via retardation waveplates.

\subsection{Reflection in CNOT basis}

Rotating the photonic basis states to

\begin{align*}
    \ket{+} &= \frac{\ket{V} + \ket{\pi}}{\sqrt{2}}\\
    \ket{-} &= \frac{\ket{V} - \ket{\pi}}{\sqrt{2}}
\end{align*}

and keeping the atom either in the coupled ($\ket{1}$) or uncoupled case ($\ket{0}$), results in the following reflection:

\begin{align*}
    \ket{0,+} &\rightarrow \ket{0,-}\\
    \ket{0,-} &\rightarrow \ket{0,+}\\
    \ket{1,+} &\rightarrow \ket{1,-}\\
    \ket{1,+} &\rightarrow \ket{1,-}\\
\end{align*}

This then tells us, that simply by using a superposition of our CPF basis, our CPF gate now represents an atom-photon controlled-NOT (CNOT) gate.

\notetome{Maybe also include brewster plate stuff here, although not quite sure}
